% Appendix C -- Procedurer

\chapter{Procedurer}
\label{AppendixC}

Dette bilag samler de praktiske procedurer, der bruges til at
idriftsætte, opstarte og fejlsøge robotten i den konfiguration, som
rapporten omfatter. Procedurerne er en opdatering af det eksisterende
materiale, som indgår i workspace-repositoriet.

% ---------------------------------------------------------------------
\section{Opstartsprocedure (opdateret)}
\label{sec:appC-opstart}

% TODO: Den opdaterede opstartsprocedure for robotten:
%  - kabel-tjekliste (forsyning, USB-bokse på korrekt port)
%  - launch-rækkefølge (servos.launch.py før slider_control osv.)
%  - validering: hvordan man bekræfter at hver arm reagerer korrekt
%  - sikkerhedshensyn (forsyning slukket før kabelændringer)

\textit{Indhold tilføjes -- baseres på den eksisterende opstartsprocedure
i workspace-repositoriet, opdateret med ændringerne fra
kapitel~\ref{ch:el} og kapitel~\ref{ch:idrift}.}

% ---------------------------------------------------------------------
\section{Portbinding -- procedure}
\label{sec:appC-portbinding}

% TODO: Trin-for-trin procedure for at finde og verificere
% /dev/serial/by-path/-stier til hver ESP32-kontroller:
%  - kommandoer (udls -l /dev/serial/by-path/, dmesg | grep tty osv.)
%  - hvordan man identificerer hvilken sti der svarer til hvilken farve
%  - hvordan manager-nodens konfiguration opdateres med de nye stier
% Selve konfigurationsfilen ligger i Bilag~\ref{AppendixB},
% afsnit~\ref{sec:appB-portbinding}.
%
% TODO ALEX: bekræft port-til-kabel-mapping mod den fysiske montage.
% Foreløbig mapping (skal verificeres):
%   Port 2.2 -- GULT kabel  -- venstre arm
%   Port 2.3 -- RØDT kabel  -- højre arm og hoved
%   Port 2.1 -- HVIDT kabel -- hænder

\textit{Indhold tilføjes.}

% ---------------------------------------------------------------------
\section{Kalibrerings- og idriftsættelsesprocedure}
\label{sec:appC-kalibrering}

% TODO: Procedure for at sætte hver servo i defaultposition og
% kontrollere software-grænser (jf.~Tabel~\ref{tab:kalibrering-tab}
% i Bilag~\ref{AppendixA}).

\textit{Indhold tilføjes.}

% ---------------------------------------------------------------------
\section{Kendte fejlbilleder og afhjælpning}
\label{sec:appC-fejlbilleder}

% TODO: Liste over de fejlbilleder, der er observeret under
% idriftsættelsen, og hvordan de afhjælpes:
%  - "armen kører til forkert position" -> tjek by-path-mapping
%  - "armen oscillerer" -> feedback_enabled må ikke stå til true
%  - "ESP32 svarer ikke" -> tjek USB-forbindelse og strømforsyning

\textit{Indhold tilføjes.}

% ---------------------------------------------------------------------
\section{Opdaterede README-filer}
\label{sec:appC-readme}

% TODO: De README-filer der er ændret/opdateret i workspace-pakkerne
% som del af projektet. Indsæt enten som verbatim-blokke eller som
% kort referat med markering af hvad der er ændret.

\textit{Indhold tilføjes.}
